% Source: physical PDF pp. 258--263 (printed pp. 235--240), from the
% Section 28.7 heading through the final paragraph immediately before
% Problems.
% Inventory: Eqs. (28.7.1)--(28.7.7), three unnumbered display groups,
% Figure 28.7, twelve linked bracketed citation-marker occurrences (with
% [36] used twice), one explicit linked textual reference to [42], and no
% footnotes, tables, or centered three-asterisk dividers.
% Corrected the opening paragraph's equation references. Uncertainties: none.

\subsection{Baryon and Lepton Non-Conservation}

The extra particles in supersymmetric models provide several new mechanisms
for baryon and lepton non-conservation. We saw in Section 28.1 that there are
various baryon- and lepton-number non-conserving supersymmetric operators
\eqref{eq:28.1.4} and \eqref{eq:28.1.5} of dimensionality four that can be
included in a renormalizable
\(SU(3)\times SU(2)\times U(1)\)-invariant theory and that
would lead to processes like proton decay at a catastrophic rate. These
terms can be excluded from the Lagrangian by imposing \(R\) parity
conservation (or, equivalently, invariance under a change of sign of all
quark and lepton chiral superfields) but this does not exclude various
\(SU(3)\times SU(2)\times U(1)\)-invariant but baryon- and
lepton-number non-conserving operators of dimensionality \(d>4\). As
discussed in Section 21.3, if there is an underlying mechanism for baryon
and lepton non-conservation characterized by some high mass scale \(M\),
then these operators will appear in the effective Lagrangian of the
standard model with coefficients proportional to \(M^{4-\dd}\). With only
the fields of the non-supersymmetric standard model, operators that can
violate baryon conservation have a minimum dimensionality
six~\hyperref[ch28-ref-36]{[36]}, and thus give baryon
non-conserving amplitudes proportional to \(M^{-2}\). The new fields
required by supersymmetry lead to two important changes in estimates of
baryon non-conserving processes like proton and bound neutron decay. As we
saw in Section 28.2, the change in the renormalization group equations gives
larger estimates of \(M\), which decreases the effect of the dimension six
operators. At the same time, these new fields allow the construction of new
operators of dimensionality five, which give baryon non-conserving
amplitudes proportional to \(M^{-1}\) and are therefore likely to make the
dominant contribution to proton and bound neutron
decay~\hyperref[ch28-ref-37]{[37]}.

The supersymmetric operators of dimensionality five that can be formed out
of chiral superfields (generically called \(\Phi\)) are of the forms
\[
  (\Phi^*\Phi\Phi)_D
  \quad\hbox{and}\quad
  (\Phi\Phi\Phi\Phi)_{\mathcal F}\,,
\]
and their complex conjugates. (We do not consider operators that include
derivatives or gauge fields, because they do not offer any additional
possibility of baryon or lepton non-conservation.)

\IfFileExists{chapters/chapter28/figures/fig28-07.tex}{%
  % Source: physical PDF p. 259 (printed p. 236). Coverage: complete Figure
% 28.7 and caption. The topology, line styles, arrows, and central dot were
% reconstructed from the full-resolution rendered scan.
\begin{figure}[H]
  \centering
  \begin{tikzpicture}[
    x=1cm,
    y=1cm,
    line cap=round,
    line join=round,
    fermion/.style={
      draw,
      line width=0.85pt,
      postaction={decorate},
      decoration={
        markings,
        mark=at position 0.56 with {\arrow{Stealth[length=2.0mm,width=1.6mm]}}
      }
    },
    scalar/.style={
      draw,
      line width=0.85pt,
      dash pattern=on 3.2pt off 3.2pt
    },
    gaugino/.style={
      draw,
      line width=0.8pt,
      decorate,
      decoration={snake,amplitude=1.15mm,segment length=4.8mm}
    }
  ]
    \coordinate (fvertex) at (-0.55,0);
    \coordinate (upper) at (2.05,1.20);
    \coordinate (lower) at (2.05,-1.20);

    \draw[fermion] (-3.25,1.75) -- (fvertex);
    \draw[fermion] (-3.25,-1.75) -- (fvertex);
    \fill (fvertex) circle[radius=0.10];

    \draw[scalar] (fvertex) -- (upper);
    \draw[scalar] (fvertex) -- (lower);

    \draw[gaugino] (upper) -- (lower);
    \draw[line width=0.65pt] (upper) -- (lower);

    % All four external fermion lines run into the diagram, as in the source.
    \draw[fermion] (4.25,1.85) -- (upper);
    \draw[fermion] (4.25,-1.85) -- (lower);
  \end{tikzpicture}
  \renewcommand{\thefigure}{28.7}
  \caption{A diagram that can produce a four-fermion interaction among
  quarks and/or leptons that violates baryon and lepton number conservation.
  Here solid lines are quarks and/or leptons; dashed lines are squarks and/or
  sleptons; the combined solid and wavy line is a gaugino; and the dot is a
  vertex arising directly from the \(\mathcal F\)-term interactions
  \eqref{eq:28.7.3}.}
  \label{fig:28.7}
\end{figure}
%
}{%
  % Source: physical PDF p. 259 (printed p. 236). Coverage: complete Figure
% 28.7 and caption. The topology, line styles, arrows, and central dot were
% reconstructed from the full-resolution rendered scan.
\begin{figure}[H]
  \centering
  \begin{tikzpicture}[
    x=1cm,
    y=1cm,
    line cap=round,
    line join=round,
    fermion/.style={
      draw,
      line width=0.85pt,
      postaction={decorate},
      decoration={
        markings,
        mark=at position 0.56 with {\arrow{Stealth[length=2.0mm,width=1.6mm]}}
      }
    },
    scalar/.style={
      draw,
      line width=0.85pt,
      dash pattern=on 3.2pt off 3.2pt
    },
    gaugino/.style={
      draw,
      line width=0.8pt,
      decorate,
      decoration={snake,amplitude=1.15mm,segment length=4.8mm}
    }
  ]
    \coordinate (fvertex) at (-0.55,0);
    \coordinate (upper) at (2.05,1.20);
    \coordinate (lower) at (2.05,-1.20);

    \draw[fermion] (-3.25,1.75) -- (fvertex);
    \draw[fermion] (-3.25,-1.75) -- (fvertex);
    \fill (fvertex) circle[radius=0.10];

    \draw[scalar] (fvertex) -- (upper);
    \draw[scalar] (fvertex) -- (lower);

    \draw[gaugino] (upper) -- (lower);
    \draw[line width=0.65pt] (upper) -- (lower);

    % All four external fermion lines run into the diagram, as in the source.
    \draw[fermion] (4.25,1.85) -- (upper);
    \draw[fermion] (4.25,-1.85) -- (lower);
  \end{tikzpicture}
  \renewcommand{\thefigure}{28.7}
  \caption{A diagram that can produce a four-fermion interaction among
  quarks and/or leptons that violates baryon and lepton number conservation.
  Here solid lines are quarks and/or leptons; dashed lines are squarks and/or
  sleptons; the combined solid and wavy line is a gaugino; and the dot is a
  vertex arising directly from the \(\mathcal F\)-term interactions
  \eqref{eq:28.7.3}.}
  \label{fig:28.7}
\end{figure}
%
}

In the notation of Section 28.1, the
\(SU(3)\times SU(2)\times U(1)\)-invariant operators of
dimensionality five that also conserve \(R\) parity are
\begin{equation}
  (LLH_2H_2)_{\mathcal F}\,,
  \tag{28.7.1}\label{eq:28.7.1}
\end{equation}
\begin{equation}
  (L\bar E H_2^*)_D\,,
  \qquad
  (Q\bar D H_2^*)_D\,,
  \qquad
  (Q\bar U H_1^*)_D\,,
  \qquad
  (QQ\bar U\bar D)_{\mathcal F}\,,
  \qquad
  (Q\bar U L\bar E)_{\mathcal F}\,,
  \tag{28.7.2}\label{eq:28.7.2}
\end{equation}
and
\begin{equation}
  (QQQL)_{\mathcal F}\,,
  \qquad
  (\bar U\bar U\bar D\bar E)_{\mathcal F}\,,
  \tag{28.7.3}\label{eq:28.7.3}
\end{equation}
with obvious contraction of indices as dictated by \(SU(3)\) and \(SU(2)\)
conservation. The interaction \eqref{eq:28.7.1} is the supersymmetric
version of the dimensionality five operator which would provide small
neutrino masses in some theories~\hyperref[ch28-ref-38]{[38]}. The interactions
\eqref{eq:28.7.2} only provide small corrections to processes that already
occur in the renormalizable terms of the supersymmetric standard model. It
is the interactions \eqref{eq:28.7.3} that provide new mechanisms for
baryon as well as lepton non-conservation.

According to Eq.~\eqref{eq:26.4.4}, the quarks and leptons enter in the
interactions \eqref{eq:28.7.3} through terms involving a pair of quark
and/or lepton fields and a pair of squark and/or slepton fields. In order to
generate reactions among quarks and leptons alone, it is necessary for the
pair of squarks and/or sleptons to be converted into a pair of quarks
and/or leptons by exchanging a gaugino in the one-loop diagram shown in
Figure~\ref{fig:28.7}. This will produce effective four-fermion \(qqq\ell\)
interactions with \(d=6\) among three quarks and a lepton. The coupling
\(g_6\) of these interactions will be proportional to the square of the
gauge coupling \(g\) or \(g'\) or \(g_s\) of the gaugino, to the
supersymmetry-breaking mass of the gaugino, to the inverse square of the
larger of the gaugino and squark or slepton masses (which is needed to give
a coupling of the right dimensionality), and to a factor of order
\(1/8\pi^2\) from the loop integration.

It might be thought that the greater strength of the gluino coupling would
make gluino exchange the dominant contribution to \(g_6\). (Indeed, in
gauge-mediated supersymmetry-breaking theories, for moderate values of the
trace \eqref{eq:28.6.3} and with \(g\approx g'\),
Eqs.~\eqref{eq:28.6.13}--\eqref{eq:28.6.14} give
\[
  m_{\mathrm{gluino}}\approx m_{\mathrm{squark}}
    \approx\frac{g_s^2}{16\pi^2}M_*\,,
  \qquad
  m_{\mathrm{wino}}\approx m_{\mathrm{slepton}}\approx m_{\mathrm{bino}}
    \approx\frac{g^2}{16\pi^2}M_*\,,
\]
where \(M_*\) is a mass characterizing the messenger sector. Therefore in
such theories individual gluino exchange diagrams give contributions
proportional to \(g_s^2/m_{\mathrm{gluino}}\), while wino or bino exchange
(or, more precisely, chargino or neutralino exchange) makes a contribution
proportional to \(g^2m_{\mathrm{wino}}/m_{\mathrm{squark}}^2\), which is
smaller by a factor roughly equal to
\(m_{\mathrm{wino}}g_s^2/m_{\mathrm{gluino}}g^2
\approx g^4/g_s^4\).) However, there is a cancellation among the diagrams
for gluino exchange that strongly suppresses their contribution. This was
originally shown using a Fierz identity among four-fermion
operators~\hyperref[ch28-ref-39]{[39]}, but the same result
can be obtained with no equations at all. To conserve color, the
coefficients of the operators \((QQQL)_{\mathcal F}\) and
\((\bar U\bar U\bar D\bar E)_{\mathcal F}\) must be totally antisymmetric
in the colors of the three quark or antiquark superfields, and since these
superfields are bosonic, they must then also be antisymmetric in their
flavors as well. Gluino interactions are flavor-independent, so if we can
neglect the flavor dependence of the squark masses then the coefficients of
the four-fermion \(d=6\) operators will also be totally antisymmetric in
flavor as well as color. Fermi statistics then requires the coefficients
of these operators to be also totally antisymmetric in the spin indices of
the quark or antiquark fields. But the three quark or antiquark fields in
the \(d=6\) operators derived by gaugino exchange from
\((QQQL)_{\mathcal F}\) or \((\bar U\bar U\bar D\bar E)_{\mathcal F}\)
are all left-handed and therefore have only two independent spin indices,
so no coefficient can be antisymmetric in all three spins. The contribution
of gluino exchange to the \(d=6\) operators would therefore vanish if the
squark masses were all equal and so this contribution is suppressed by the
fractional differences among the masses of the different squarks. According
to Eq.~\eqref{eq:28.6.4}, in theories of gauge-mediated supersymmetry
breaking the fractional differences between the \(\bar U\) and \(\bar D\)
squark masses are of order \(g'^4/g_s^4\), so gluino exchange between the
antisquarks in the \((\bar U\bar U\bar D\bar E)_{\mathcal F}\) operator
generates a dimension six four-fermion operator with a coefficient of the
same order as that generated by bino exchange. However, since gluinos
conserve flavor this operator like the
\((\bar U\bar U\bar D\bar E)_{\mathcal F}\) operator must be totally
antisymmetric in antiquark flavors, so that it must involve \(c\) or \(t\)
quarks and therefore cannot contribute directly to proton or bound neutron
decay. On the other hand, Eq.~\eqref{eq:28.6.4} indicates that the
fractional mass differences among the \(Q\) squarks of different flavors is
much less than of order \(g^4/g_s^4\), so that gluino exchange between the
squarks in the \((QQQL)_{\mathcal F}\) operator makes a contribution to
\(g_5\) that is much smaller than wino or bino exchange. We conclude that
at least in the theories of gauge-mediated supersymmetry breaking, gluino
exchange makes a smaller contribution to proton or bound neutron decay than
wino or bino exchange. In other models gluino exchange may make a
contribution comparable to that of other
processes~\hyperref[ch28-ref-40]{[40]}.

For \(g\approx g'\) and \(m_{\mathrm{wino}}\approx m_{\mathrm{bino}}\),
the contribution of wino or bino exchange to the dimension six operators
is of order
\begin{equation}
  g_6\approx
  \frac{g^2g_5m_{\mathrm{wino}}}
       {8\pi^2m_{\mathrm{squark}}^2}\,,
  \tag{28.7.4}\label{eq:28.7.4}
\end{equation}
where \(g_5\) is a typical value of the couplings of the effective \(d=5\)
interactions \eqref{eq:28.7.3}. If the wino and squark masses have the same
ratio \((g^2/g_s^2)\) as in theories of gauge-mediated supersymmetry
breaking, then this gives
\begin{equation}
  g_6\approx
  \frac{g^4g_5}
       {8\pi^2g_s^2m_{\mathrm{squark}}}\,.
  \tag{28.7.5}\label{eq:28.7.5}
\end{equation}

The four-fermion \(qqq\ell\)-terms of dimensionality six in the effective
Lagrangian are the same as those that had been supposed to generate
processes like proton decay in non-supersymmetric
theories~\hyperref[ch28-ref-36]{[36]}. They produce proton
and bound neutron decay at a rate which on dimensional grounds must be of
the form
\begin{equation}
  \Gamma_N=c_Nm_N^5g_6^2\,,
  \tag{28.7.6}\label{eq:28.7.6}
\end{equation}
where \(c_N\) is a pure number that must be calculated by non-perturbative
calculations in quantum chromodynamics. Much work has gone into these
calculations, with results~\hyperref[ch28-ref-41]{[41]} generally in the range
\(c_N\approx 3\times10^{-3\pm0.7}\).

To estimate \(g_5\), we note that it is not possible to produce
\(\mathcal F\)-terms like \eqref{eq:28.7.3} that involve only left-chiral
superfields by the tree-approximation exchange of gauge supermultiplets,
which always interact with both left-chiral superfields and their
right-chiral complex conjugates. Thus the interactions \eqref{eq:28.7.3}
arise in the tree approximation only from the exchange of particles of
chiral superfields, and therefore \(g_5\) is of the order of \(g_T^2/M_T\),
where \(g_T\) is a typical baryon and lepton non-conserving coupling of some
superheavy left-chiral superfield of mass \(M_T\) to the quark and lepton
superfields. To produce the interactions \eqref{eq:28.7.3}, these
superheavy particles must be color triplets or antitriplets, and \(SU(2)\)
triplets or singlets. Whatever gauge group unifies the strong and
electroweak interactions presumably dictates some relation between the
interactions of the superheavy color triplet \(T\) and the familiar color
singlet \(H_1\) and \(H_2\). Then \(g_T\) will be of the same order as the
Yukawa coupling constants in the interactions \eqref{eq:28.1.2} and
\eqref{eq:28.1.3} that give mass to the known quarks and leptons, and which
are equal to quark or lepton masses divided by the vacuum expectation
values of order \(G_F^{-1/2}\simeq300\ \mathrm{GeV}\) of
\(\mathcal H_1^0\) or \(\mathcal H_2^0\). We therefore take
\begin{equation}
  g_5\approx\frac{G_Fm_f^2}{M_T}\,,
  \tag{28.7.7}\label{eq:28.7.7}
\end{equation}
where \(m_f\) is some typical quark or lepton mass. We have seen that the
dimension five operators are antisymmetric in quark flavors, so as a
compromise between the masses of the \(s\) quark and the \(u\) or \(d\)
quarks, we will take \(m_f=30\ \mathrm{MeV}\). Combining
Eqs.~\eqref{eq:28.7.5}--\eqref{eq:28.7.7} and taking
\(M_T=2\times10^{16}\ \mathrm{GeV}\) (as suggested by the results of
Section 28.2), \(c_N=0.003\), \(g_s^2/4\pi=0.118\),
\(g^2/4\pi=1/(0.23\times137)\), and
\(m_{\mathrm{squark}}=1\ \mathrm{TeV}\), we find a proton (or bound
neutron) lifetime \(\Gamma_N^{-1}\) of about \(2\times10^{31}\)
years~\hyperref[ch28-ref-42]{[42]}. This is not very
different from experimental lower bounds on the partial lifetimes of what
are expected to be the leading modes of proton decay, variously quoted as
ranging from \(10^{31}\) to \(5\times10^{32}\) years. At the time of
writing the most stringent bounds are set by the non-observation of proton
decay in the large Super Kamiokande neutrino detector in
Japan~\hyperref[ch28-ref-42a]{[42a]}: the partial lifetimes
for the decays \(p\longrightarrow e^+\pi^0\) and
\(p\longrightarrow\bar\nu K^+\) are greater than
\(2.1\times10^{33}\) years and \(5.5\times10^{32}\) years, respectively.
There is an uncertainty of a factor of at least 100 in the above estimate
of the theoretical lifetime from the uncertainty in the squark mass alone,
so it is too soon to say that there is any discrepancy between experiment
and theoretical expectations. On the other hand, supersymmetry raises the
possibility that baryon non-conservation may be discovered soon.

We can also say a little on general grounds about the expected branching
ratios for various proton and bound neutron decay modes. As we have
mentioned, the dimension five operators \eqref{eq:28.7.3} must be totally
antisymmetric in the flavor of the quark superfields, so the only operators
that concern us here are of the forms
\((U_iD_jD_kN_\ell)_{\mathcal F}\),
\((D_iU_jU_kE_\ell)_{\mathcal F}\), and
\((\bar D_i\bar U_j\bar U_k\bar E_\ell)_{\mathcal F}\), where
\(i,j,k,\ell\) are generation indices, and in each case \(j\ne k\). The
exchange of neutral winos or binos then produces \(d=6\) four-fermion
operators of the forms \(u_id_jd_k\nu_\ell\), \(d_iu_ju_ke_\ell\), and
\(\bar d_i\bar u_j\bar u_k\bar e_\ell\), with \(j\ne k\) and \(i\) and
\(\ell\) arbitrary, while charged wino exchange produces the same
four-fermion operators with \(i\ne j\) and \(k\) and \(\ell\) arbitrary.
The only quarks light enough to be involved in proton decay are \(u\),
\(s\), and \(d\); ignoring all the others and the small mixing angles in
the third generation, we have
\[
  u_1=u\,,
  \qquad
  d_1=\dd\cos\theta_c+s\sin\theta_c\,,
  \qquad
  d_2=-\dd\sin\theta_c+s\cos\theta_c\,,
\]
where \(\theta_c\) is the Cabibbo angle, while \(u_2\), \(u_3\), and \(d_3\)
can be ignored. The four-fermion operators that can be produced by wino or
bino exchange and can contribute to proton or bound neutron decay thus are
\(uds\nu_\ell\cos(2\theta_c)\), \(udd\nu_\ell\sin(2\theta_c)\),
\(uuse_\ell\cos\theta_c\), and \(uude_\ell\sin\theta_c\), plus others with
quarks and leptons replaced with antiquarks and antileptons. All other
things being equal, the dominant decay modes are therefore
\(p\longrightarrow K^+\bar\nu\), \(n\longrightarrow K^0\bar\nu\),
\(p\longrightarrow K^0e^+\), and \(p\longrightarrow K^0\mu^+\), while the
rates for the decay modes \(p\longrightarrow\pi^+\bar\nu\),
\(n\longrightarrow\pi^0\bar\nu\), \(p\longrightarrow\pi^0e^+\),
\(p\longrightarrow\pi^0\mu^+\), and \(n\longrightarrow\pi^-e^+\) are
suppressed by factors \(\sin^2\theta_c=0.05\), though also enhanced
somewhat by the greater phase space available.

These considerations do not lead to definite predictions of branching
ratios, because in addition to all the factors mentioned above, the
coefficients generically called \(g_5\) of the operators
\eqref{eq:28.7.3} may have a strong dependence on the superfield flavors
appearing in these operators. To go further, one needs a specific theory
for the generation of the dimension five operators. Most of the authors of
Reference~\hyperref[ch28-ref-42]{42} concluded on the basis of a
supersymmetric version of \(SU(5)\) theories that proton and bound neutron
decay would be dominated by the processes \(p\longrightarrow K^+\bar\nu\)
and \(n\longrightarrow K^0\bar\nu\), but for a model based on \(SO(10)\)
charged lepton modes can become prominent~\hyperref[ch28-ref-43]{[43]}.
Also, in some models higgsino exchange can compete with wino and bino
exchange~\hyperref[ch28-ref-44]{[44]}, increasing the rate
of \(p\longrightarrow K^+\bar\nu\). It seems a good idea in searches for
baryon non-conservation to keep an open mind as to the decay modes to be
expected in proton or bound neutron decay.

Of course, it is possible that all of these baryon non-conserving processes
are prohibited by some sort of conservation law. As mentioned in Section
28.1, string theory argues against baryon conservation being a fundamental
global continuous symmetry, but the baryon non-conserving operators
\eqref{eq:28.7.3} may be forbidden by a \(\mathbb Z_3\) multiplicative
symmetry known as baryon parity~\hyperref[ch28-ref-45]{[45]}, under which the \(Q\)
superfield is neutral; the \(H_2\) and \(\bar D\) superfields are multiplied
by the phase \(\exp(\ii\pi/3)\), and the \(L\), \(H_1\), \(\bar U\), and
\(\bar E\) superfields are multiplied by the opposite phase
\(\exp(-\ii\pi/3)\). This symmetry is designed to allow the fundamental
Yukawa couplings \eqref{eq:28.1.2} and \eqref{eq:28.1.3} as well as the
\(\mu\)-term \eqref{eq:28.5.7} and the lepton non-conserving terms
\eqref{eq:28.1.4} and \eqref{eq:28.7.1}, but it rules out the dimension
four baryon non-conserving term \eqref{eq:28.1.5} and the dimension five
baryon non-conserving terms \eqref{eq:28.7.3}. This symmetry is
spontaneously broken by the appearance of vacuum expectation values of
\(\mathcal H_1^0\) and \(\mathcal H_2^0\) (and perhaps the sneutrino fields
\(\mathcal N\)), and with no conservation law for \(R\) parity, there would
be nothing to keep the lightest supersymmetric particle stable.
